\chapter{Architecture Overview}\label{ch:architecture}

Part~II looks inside the code generator. This chapter gives the big picture: the
design principles, the package structure, and the driver that runs the compilation
pipeline. The following chapters then look at the pieces: the intermediate
representation (\cref{ch:ir}), the translation from SymPy (\cref{ch:translation}), the
transformations (\cref{ch:transformations}), and the hardware platforms
(\cref{ch:platforms,ch:rvv-backend}).

\section{Design principles}

\begin{description}
  \item[Two languages.] Users write kernels in a \emph{mathematical} language: SymPy
        expressions over fields, mostly untyped, with no notion of loops or memory. The
        backend works on an \emph{imperative, typed} intermediate representation (IR): an
        abstract syntax tree of C-like statements and expressions in which every node has
        a data type. The translation between the two (``freeze'' and ``typify'')
        happens once, at the beginning.
  \item[Small, composable passes.] Every step of the compiler is a class whose
        instances are called on an AST and return a new or modified AST:
        \code{EliminateConstants(ctx)(ast)}, \code{LowerToC(ctx)(ast)}, \dots. The
        driver is a short, readable sequence of such calls, and a developer can run any
        subset by hand (\cref{sec:manual-pipeline}).
  \item[A shared context.] All passes of one kernel share one
        \cls{KernelCreationContext}. It owns the symbol table, the fields and their
        memory buffers, the iteration space, the reductions and the required headers.
  \item[Iteration strategies as data.] How the iteration space is traversed (loops,
        OpenMP, SIMD blocks with remainder loops, GPU thread indices) is first expressed
        as \emph{iteration axes} inside the AST. It is lowered to concrete loops only
        afterwards. CPU and GPU kernels share this machinery and differ in the
        strategy they choose (\cref{sec:axes}).
  \item[Platforms encapsulate the hardware.] Everything that depends on the instruction
        set (vector intrinsics, math function names, reductions, required headers) lives
        in a \cls{Platform} class (\cref{ch:platforms}). A new architecture is added by
        writing a new platform, not by changing the passes (\cref{sec:new-backend}).
  \item[Self-describing results.] The output is a \code{Kernel} object with typed
        parameters and semantic properties (\cref{sec:kernel-object}). Compiling it is
        a separate concern, handled by the JIT or by external code generators.
\end{description}

\section{Package structure}

\Cref{tab:packages} maps the packages under \file{src/pystencils/} to these roles.
\Cref{app:source-map} gives a file-by-file map.

\begin{table}[htbp]
  \centering\small
  \caption{Main packages of \ps{} 2.0.}\label{tab:packages}
  \begin{tabularx}{\linewidth}{@{}l >{\raggedright\arraybackslash}p{4.4cm} >{\raggedright\arraybackslash}X@{}}
    \toprule
    \textbf{Layer} & \textbf{Package / module} & \textbf{Responsibility}\\
    \midrule
    \multirow{5}{*}{Frontend}
      & \file{field.py} & \code{Field}, \code{Field.Access}, layouts\\
      & \file{assignment.py}, \file{simp/} & assignments, collections, simplification\\
      & \file{sympyextensions/} & typed symbols, casts, reductions, RNG, integer functions\\
      & \file{fd/} & finite-difference discretization\\
      & \file{flow/} & experimental flowgraph language\\
    \midrule
    \multirow{2}{*}{Codegen API}
      & \file{codegen/} & \code{create_kernel}, driver, config, \code{Target}, \code{Kernel}, \code{Parameter}, GPU launch configurations\\
      & \file{types/} & the type system (\code{PsType} hierarchy, \code{create_type})\\
    \midrule
    \multirow{5}{*}{Backend}
      & \file{backend/kernelcreation/} & context, analysis, iteration spaces, freeze, typify, AST factory\\
      & \file{backend/ast/} & IR node classes, traversal and analysis helpers\\
      & \file{backend/transformations/} & optimization and lowering passes\\
      & \file{backend/platforms/} & hardware platforms and intrinsic selection\\
      & \file{backend/emission/} & C and IR printers\\
    \midrule
    \multirow{2}{*}{Runtime}
      & \file{jit/} & \code{CpuJit}, \code{CupyJit}, \code{SYCLJit}, \code{no_jit}\\
      & \file{include/} & C++ runtime headers (SIMD helpers, Philox, GPU support)\\
    \bottomrule
  \end{tabularx}
\end{table}

\section{The kernel creation driver}\label{sec:driver}

\code{create_kernel} creates a \cls{DefaultKernelCreationDriver}
(\file{codegen/driver.py}) and calls it. The constructor
\begin{enumerate}
  \item checks that at most one of \code{ghost_layers}, \code{iteration_slice} and
        \code{index_field} is set (and falls back to \code{ghost_layers=AUTO});
  \item creates the \cls{KernelCreationContext} with the configured \code{default_dtype}
        and \code{index_dtype};
  \item resolves automatic targets (\code{CurrentCPU} becomes \code{auto_cpu()}), and for
        GPU targets creates the GPU indexing object that later provides the launch
        configuration;
  \item selects the \cls{Platform} for the target (\cref{tab:platform-selection}).
\end{enumerate}
Its \code{__call__} method then runs the five stages of \cref{fig:pipeline}:
\begin{pycode}
def __call__(self, symbolic_kernel: SymbolicKernel) -> Kernel:
    kernel_body = self.parse_kernel_body(symbolic_kernel)
    kernel_ast = self._materialize_iteration_space(kernel_body)
    kernel_ast = self._general_optimize(kernel_ast)
    kernel_ast = self._lowering(kernel_ast)
    return self._finalize(kernel_ast)
\end{pycode}

\begin{figure}[htbp]
  \centering
  % The five stages of DefaultKernelCreationDriver.__call__
\begin{tikzpicture}[
  font=\small,
  stage/.style={draw=psblue!80!black,fill=psblue!8,rounded corners=2pt,
    minimum width=3.4cm,minimum height=1.0cm,align=center,inner sep=3pt},
  passes/.style={draw=psgray!50,fill=white,rounded corners=2pt,align=left,
    text width=8.6cm,inner sep=4pt,font=\footnotesize},
  io/.style={draw=psteal!80!black,fill=psteal!8,rounded corners=2pt,align=center,
    inner sep=4pt,font=\footnotesize},
  arr/.style={-{Stealth[length=2.2mm]},thick,draw=psgray},
  conn/.style={draw=psgray!60,dashed},
]
  \node[io] (in) {SymPy assignments / \texttt{AssignmentCollection} / flowgraph\\
    \texttt{CreateKernelConfig}};

  \node[stage,below=0.55cm of in] (s1) {\textbf{1. parse}\\\scriptsize\texttt{parse\_kernel\_body}};
  \node[stage,below=0.95cm of s1] (s2) {\textbf{2. iteration space}\\\scriptsize\texttt{\_materialize\_iteration\_space}};
  \node[stage,below=0.95cm of s2] (s3) {\textbf{3. optimize}\\\scriptsize\texttt{\_general\_optimize}};
  \node[stage,below=0.6cm of s3] (s4) {\textbf{4. lower}\\\scriptsize\texttt{\_lowering}};
  \node[stage,below=0.6cm of s4] (s5) {\textbf{5. finalize}\\\scriptsize\texttt{\_finalize}};
  \node[io,below=0.55cm of s5] (out) {\texttt{Kernel} / \texttt{GpuKernel}};

  \foreach \a/\b in {in/s1,s1/s2,s2/s3,s3/s4,s4/s5,s5/out} \draw[arr] (\a) -- (\b);

  \node[passes,right=0.6cm of s1] (p1) {\texttt{KernelAnalysis} (checks, field registry)
    $\to$ create iteration space (ghost layers / slice / index field)
    $\to$ \texttt{FreezeExpressions} (SymPy $\to$ IR) $\to$ \texttt{Typifier}};
  \node[passes,right=0.6cm of s2] (p2) {CPU: \texttt{DefaultCpuLoopStrategies} / GPU:
    \texttt{GpuIndexMappingStrategy} build an \emph{axes tree}
    (\texttt{AxisExpansion}, \texttt{HoistIterationInvariantDeclarations})
    $\to$ \texttt{MaterializeAxes} (loops, OpenMP, SIMD, GPU indices) $\to$
    \texttt{ReductionsToMemory}};
  \node[passes,right=0.6cm of s3] (p3) {\texttt{CanonicalizeSymbols} $\to$
    \texttt{EliminateConstants} (folding, constant extraction) $\to$
    \texttt{HoistIterationInvariantDeclarations}};
  \node[passes,right=0.6cm of s4] (p4) {vector CPU: platform's \texttt{SelectIntrinsics}
    $\to$ \texttt{LowerToC} (linearize buffer accesses) $\to$
    \texttt{SelectFunctions} (platform-specific math)};
  \node[passes,right=0.6cm of s5] (p5) {\texttt{CanonicalizeSymbols} $\to$
    \texttt{KernelFactory}: collect parameters and headers, attach JIT and launch
    configuration};

  \foreach \i in {1,...,5} \draw[conn] (s\i.east) -- (p\i.west);
\end{tikzpicture}

  \caption[The five stages of the kernel creation driver]{The five stages of \texttt{DefaultKernelCreationDriver} and the passes they run.}
  \label{fig:pipeline}
\end{figure}

\paragraph{Stage 1: parse.}
For assignment input, \code{parse_assignment_collection} wraps the input in an
\code{AssignmentCollection}, runs \cls{KernelAnalysis} (constraint checks; registers all
fields in the context and marks read-only fields \code{const}), creates the iteration
space, and translates the assignments into a typed \code{PsBlock} with
\cls{FreezeExpressions} and \cls{Typifier} (\cref{ch:translation}). Flowgraph input
(\cref{sec:flow}) takes a separate path, \code{parse_flowgraph}: it registers the fields
directly and uses \cls{FreezeFlowgraph}, which typifies each node on its own.
\cls{KernelAnalysis} is not run on this path.

\paragraph{Stage 2: materialize the iteration space.}
For CPU and GPU platforms, a strategy object builds an \emph{axes tree} from the
iteration space. For CPUs this is \cls{DefaultCpuLoopStrategies}
(\file{codegen/cpu_loop_strategies.py}), which reads the OpenMP, vectorization and
blocking options. For GPUs it is \cls{GpuIndexMappingStrategy}
(\file{codegen/gpu_indexing.py}), which reads the indexing scheme. Then
\cls{MaterializeAxes} turns the axes into loops, pragmas, vectorized bodies or GPU index
computations, and \cls{ReductionsToMemory} adds the initialization and write-back of
reduction variables (\cref{ch:transformations}). The SYCL platform still uses the older
platform hook \code{materialize_iteration_space} instead.

\paragraph{Stage 3: general optimizations.}
\cls{CanonicalizeSymbols} renames duplicate declarations and makes read-only symbols
\code{const}. \cls{EliminateConstants} folds constant expressions and extracts constant
subexpressions into \code{__c_*} declarations. \cls{HoistIterationInvariantDeclarations}
moves loop-invariant declarations out of loops. No option controls this stage.

\paragraph{Stage 4: lowering.}
For vector CPU platforms, the platform's \cls{SelectIntrinsics} pass replaces vector
operations with intrinsic calls. \cls{LowerToC} linearizes multidimensional buffer
accesses into pointer arithmetic, and \cls{SelectFunctions} asks the platform to replace
abstract math functions (\code{exp}, \code{min}, $\pi$, \dots) with concrete C
functions.

\paragraph{Stage 5: finalize.}
A last \cls{CanonicalizeSymbols} pass handles symbols introduced during lowering. Then
\cls{KernelFactory} collects all symbols that are used but never defined; these become
the parameters. It translates internal buffer properties into \code{FieldBasePtr} and
collects headers from the types, the platform and the context. It also attaches the JIT,
the context's metadata and, for GPU kernels, the launch configuration factory.

\begin{table}[htbp]
  \centering\small
  \caption{Platform selection in \texttt{DefaultKernelCreationDriver.\_get\_platform}.}
  \label{tab:platform-selection}
  \begin{tabular}{@{}ll@{}}
    \toprule
    \textbf{Target} & \textbf{Platform class} (\file{backend/platforms/})\\
    \midrule
    \code{GenericCPU} & \cls{GenericCpu}\\
    \code{X86_SSE}, \code{X86_AVX}, \code{X86_AVX512}(\code{_FP16}) & \cls{X86VectorCpu} with \code{X86VectorArch.<SSE|AVX|AVX512|...>}\\
    \code{ARM_NEON}, \code{ARM_NEON_FP16} & \cls{NeonCpu}\\
    \code{ARM_SVE} & \cls{SveCpu}\\
    \code{RISCV_RVV} & \cls{RvvCpu} (fork only)\\
    \code{CUDA}, \code{HIP} & \cls{CudaPlatform}, \cls{HipPlatform}\\
    \code{SYCL} & \cls{SyclPlatform}\\
    \bottomrule
  \end{tabular}
\end{table}

\section{Errors}

The backend reports problems through a small set of exception classes in
\file{backend/exceptions.py}, plus \code{CodegenError} for driver-level errors:
\begin{center}\small
\begin{tabular}{@{}ll@{}}
  \toprule
  \textbf{Exception} & \textbf{Raised when}\\
  \midrule
  \code{KernelConstraintsError} & the input violates SSA or independence rules\\
  \code{FreezeError} & a SymPy construct cannot be translated\\
  \code{TypificationError} & types do not match (no implicit casts)\\
  \code{VectorizationError} & an expression cannot be vectorized\\
  \code{MaterializationError} & the platform lacks an operation (e.g.\ \code{exp} in SIMD)\\
  \code{CodegenError} & invalid combination of options (e.g.\ SIMD on \code{GenericCPU})\\
  \code{PsInternalCompilerError} & a bug in \ps{} itself\\
  \bottomrule
\end{tabular}
\end{center}
